\solchapter{chem}

\sol{chem}{1} $4.5\un{eV}\times96.5\un{(kJ/mol)/eV} = 434\un{kJ/mol}$. The Boltzmann factor reaches $10^{-10}$ when $E/\kB T = 23$, at $T\approx2300\un{K}$. At room temperature $E/\kB T = 175$ and $e^{-175}\approx10^{-76}$: even with $10^{13}$ attempts per second, a given molecule would wait vastly longer than the age of the universe. (Hydrogen gas does burn readily, but through chain reactions with much lower barriers, not by thermally breaking H$_2$.)

\sol{chem}{2} $k_2/k_1 = \exp[(E_a/R)(1/T_1 - 1/T_2)]$ with $E_a/R = 6010\un{K}$. From $300$ to $310\un{K}$: a factor of $1.9$ --- the chemist's rule that rates double every $10^\circ$, and the reason a fever speeds up your metabolism. From $273$ to $373\un{K}$: a factor of about $370$, which is why food cooks.

\sol{chem}{3} $K = e^{-\Delta G^\circ/RT} = e^{8.07}\approx3200$. For $99.9\%$ completion $K = 999$, so $\Delta G^\circ = -RT\ln999 = -17\un{kJ/mol}$. A few tens of kJ/mol, a few times $RT$ in each of several factors, separates ``hardly goes'' from ``goes to completion''.

\sol{chem}{4} $e^2/(4\pi\epsilon_0r) = 1.44\un{eV\,nm}/0.236\un{nm} = 6.1\un{eV}$. Making the ions costs $5.1 - 3.6 = 1.5\un{eV}$, so the binding energy is about $4.6\un{eV}$, $10\%$ above the measured $4.2\un{eV}$. The difference is the short-range repulsion of the ions' filled electron shells (the Pauli principle), which keeps them apart and is what sets the bond length in the first place.

\sol{chem}{5} Normalize $\psi_\pm = (\psi_A\pm\psi_B)/\sqrt{2(1\pm S)}$. Then
\[
  E_\pm = \bra{\psi_\pm}\hat H\ket{\psi_\pm} = \frac{H_{AA}+H_{BB}\pm2H_{AB}}{2(1\pm S)} = \frac{H_{AA}\pm H_{AB}}{1\pm S},
\]
using $H_{BB} = H_{AA}$ by symmetry. Measure energies from $H_{AA}$ (shifting $\hat H$ by a constant changes $H_{AB}$ by $S$ times that constant; this choice makes the formulas cleanest). Then $E_\pm = \pm H_{AB}/(1\pm S)$: with $H_{AB}<0$, the bonding orbital is lowered by $\abs{H_{AB}}/(1+S)$ and the antibonding orbital raised by $\abs{H_{AB}}/(1-S)$, which is larger since $0<S<1$.

\sol{chem}{6} $E_n = n^2h^2/8mL^2$, with $h^2/8mL^2 = 0.045\un{eV}$ for $L = 2.9\un{nm}$. The highest filled level is $n = 11$ and the lowest empty one $n = 12$, so $\Delta E = 23\times0.045 = 1.03\un{eV}$, or $\lambda = 1200\un{nm}$: infrared, nowhere near the observed $450\un{nm}$ ($2.8\un{eV}$). The box model treats the chain as a uniform wire, but its bonds alternate, single and double, so the hopping alternates, $t_1$ and $t_2$. With one $\pi$ electron per carbon, the band is exactly half-filled, and Problem~\ref{prob:dimer} showed that alternating hopping opens a gap of $2\abs{t_1-t_2}$ exactly at half-filling. For a long chain the box-model gap shrinks to zero, but the dimerization gap remains: in polyacetylene, the infinitely long version of the chain, it is about $1.5$--$2\un{eV}$. Added to carotene's finite-length spacing, it pushes the absorption into the blue, and the carrot reflects the orange that is left.

\sol{chem}{7} With $\Delta H = T_b\Delta S$: water $373\times88 = 33\un{kJ/mol}$; H$_2$S $213\times88 = 19\un{kJ/mol}$. (The measured values are $41$ and $19\un{kJ/mol}$; water breaks Trouton's rule because hydrogen bonds also order the liquid, lowering its entropy.) Without hydrogen bonding, water would be like a lighter H$_2$S and would boil far below $0^\circ\text{C}$. The extra $14$--$22\un{kJ/mol}$, at $20\un{kJ/mol}$ per bond, suggests about one hydrogen bond per molecule --- each bond is shared between two molecules, so each molecule takes part in about two. Liquid water actually has nearer $3.5$ bonds per molecule, many of them bent and weakened; the estimate catches the order of magnitude.

\sol{chem}{8} Electronic energies are $E_{\text{el}}\sim\hbar^2/m_ea^2$ for a molecule of size $a$. The bond's spring constant is $k\sim E_{\text{el}}/a^2$, so the vibrational quantum is
\[
  \hbar\omega = \hbar\sqrt{k/M}\sim\sqrt{\frac{\hbar^2}{Ma^2}E_{\text{el}}} = E_{\text{el}}\sqrt{\frac{m_e}{M}}.
\]
The rotational quantum is $\hbar^2/Ma^2 = E_{\text{el}}(m_e/M) = \hbar\omega\sqrt{m_e/M}$. For a reduced mass of $10\un{u}$, $\sqrt{m_e/M}\approx1/135$: electronic $\sim10\un{eV}$ (visible and UV), vibrational $\sim0.1\un{eV}$ (infrared), rotational $\sim1\un{meV}$ (microwave). For CO, the measured values are $8\un{eV}$, $0.27\un{eV}$, and $0.5\un{meV}$. Each level of structure is a small perturbation on the one above, which is why molecular spectra can be organized at all.

\sol{chem}{9} The zero-point energies are $\tfrac12\hbar\omega$: $1500\un{cm^{-1}}$ for C--H and $1061\un{cm^{-1}}$ for C--D, a difference of $439\un{cm^{-1}}$. At $298\un{K}$, $\kB T = 207\un{cm^{-1}}$, so the ratio of rates is $e^{439/207} = e^{2.1}\approx8$. The C--H bond starts higher up its well, so it has less far to climb.

\sol{chem}{10} We have $\ln K = -\Delta H^\circ/RT + \Delta S^\circ/R$. Treating $\Delta H^\circ$ and $\Delta S^\circ$ as constant, differentiate to get $\dd\ln K/\dd T = \Delta H^\circ/RT^2$. If $\Delta H^\circ>0$ (endothermic), $K$ grows with $T$: heat is a reactant, so supplying it pushes the reaction forward (Le Chatelier). The Haber process, $\mathrm N_2 + 3\mathrm H_2\to2\mathrm{NH_3}$, is exothermic, so the yield is best at low temperature, where the reaction is hopelessly slow. The compromise is about $450^\circ\text{C}$ with an iron catalyst, and $200$ atmospheres of pressure, which favours the side with fewer molecules (two instead of four); the ammonia is removed and the unreacted gas recycled.

\sol{chem}{11} Expand $\sqrt{1\pm\kappa} = 1\pm\tfrac12\kappa - \tfrac18\kappa^2 + \dots$. The first-order terms cancel in the sum, leaving
\[
  \tfrac12\hbar(\omega_++\omega_-) = \hbar\omega_0\big(1 - \tfrac18\kappa^2+\dots\big),\qquad\Delta E = -\tfrac18\hbar\omega_0\kappa^2 = -\frac{\hbar\omega_0}{2}\frac{\alpha^2}{(4\pi\epsilon_0)^2r^6},
\]
where $\alpha = e^2/m\omega_0^2$ is the atom's polarizability. (The three-dimensional calculation replaces $\tfrac12$ by $\tfrac34$; this is London's formula.) The attraction falls off as $1/r^6$. Classically, at zero temperature, both electrons would sit at rest, with no dipoles and no force. Quantum mechanically, their zero-point motions become correlated, so that the instantaneous dipoles tend to line up, and the energy of the correlated motion is slightly lower. This van der Waals attraction holds liquid helium together and helps geckos climb walls.

\sol{chem}{12} He$_2$: two electrons fill the bonding orbital and two the antibonding orbital. By Problem~5, the antibonding orbital is raised more than the bonding orbital is lowered, so the total energy is \emph{above} that of two separate atoms: net repulsion. (The van der Waals attraction of Problem~11 does bind a helium dimer, but by only about a millikelvin, in a molecule some $5\un{nm}$ across.) He$_2^+$ has two bonding electrons and one antibonding one: the net effect is binding, and He$_2^+$ is bound by about $2.5\un{eV}$. Bond orders: He$_2$ $(2-2)/2 = 0$; He$_2^+$ $(2-1)/2 = \tfrac12$. For N$_2$, the ten valence electrons fill $\sigma_{2s}^2\sigma_{2s}^{*2}\pi_{2p}^4\sigma_{2p}^2$: eight bonding and two antibonding, a bond order of $3$. O$_2$ has two more electrons, which go into the antibonding $\pi^*$ orbitals, giving order $(8-4)/2 = 2$. Because the two $\pi^*$ orbitals have the same energy, Hund's rule puts these two electrons in different orbitals with parallel spins: O$_2$ is paramagnetic, and liquid oxygen sticks to a magnet --- a striking prediction of the molecular orbital picture that simple pictures of electron-pair bonds miss.
